\clearpage
\VerdierPhysicalPageAnchor{0208}%
\begingroup\small\setlength{\parskip}{0pt}\setlength{\emergencystretch}{1.5em}
\setlength{\abovedisplayskip}{4pt plus 2pt minus 2pt}%
\setlength{\belowdisplayskip}{4pt plus 2pt minus 2pt}%

\VerdierSemanticAnchor{ch3d-3}
\section{Study of \texorpdfstring{$\operatorname{Ext}$}{Ext}.}

\VerdierSemanticAnchor{ch3d-3-preamble}
\noindent\textit{In this section, $\mathcal A$ denotes an abelian category.}

\VerdierSemanticAnchor{ch3d-3-1}
\subsection{Different equivalent definitions.}

\VerdierSemanticAnchor{ch3d-3-1-1}
\noindent\textbf{3.1.1.}\quad
We say that the \emph{cohomological dimension of $\mathcal A$ is less than
or equal to $n$} if, for every pair $X$ and $Y$ of objects of $\mathcal A$
and every integer $p>n$, we have
$\operatorname{Ext}_{\mathcal A}^{p}(X,Y)=0$ ((1.2.5), (1.2.9)). The
\emph{cohomological dimension of $\mathcal A$}, denoted
$\operatorname{dimcoh}(\mathcal A)$, is the least integer $n$ for which the
cohomological dimension of $\mathcal A$ is less than or equal to $n$. If
there is no integer $n$ for which the cohomological dimension of $\mathcal A$
is less than or equal to $n$, we say that the cohomological dimension of
$\mathcal A$ is infinite and write
$\operatorname{dimcoh}(\mathcal A)=+\infty$.

\VerdierSemanticAnchor{ch3d-3-1-2}
\noindent\textbf{3.1.2.}\quad
Recall (1.1.3) that $\operatorname{Qis}^{*}$ denotes the set of
quasi-isomorphisms of $K^{*}(\mathcal A)$
($*=b,-,+,{}$ ``empty''). Let $X$ be an object of $K^{*}(\mathcal A)$;
$\operatorname{Qis}^{*}/X$ (respectively,
$X\backslash\operatorname{Qis}^{*}$) denotes the category of
quasi-isomorphisms with target $X$ (respectively, with source $X$)
(Chap. II, 2.2.1). Since $\operatorname{Qis}^{*}$ is a multiplicative system
of morphisms, the category $\operatorname{Qis}^{*}/X$ (respectively,
$X\backslash\operatorname{Qis}^{*}$) is filtered on the left (respectively,
on the right) (Chap. II, 2.2.1). The category
$\operatorname{Qis}^{*}/X$ (respectively,
$X\backslash\operatorname{Qis}^{*}$) has an initial (respectively, final)
object $u\colon Y\longrightarrow X$ (respectively,
$u\colon X\longrightarrow Y$) if and only if there is an object that is
$Q^{*}$-free on the left (respectively, on the right), say $Y$, and a quasi-isomorphism
$u\colon Y\longrightarrow X$ (respectively,
$u\colon X\longrightarrow Y$), as follows immediately by using
(Chap. II, 2.3.3).

\VerdierSemanticAnchor{ch3d-3-1-3}
\medskip\noindent\textbf{Proposition 3.1.3.}\quad
\textit{Let $X$ and $Y$ be two complexes of $\mathcal A$.}

\VerdierSemanticAnchor{ch3d-3-1-3-1}
\[
\textit{(1)}\qquad
\operatorname{Ext}_{\mathcal A}^{0}(X,Y)=
\varinjlim_{(\operatorname{Qis}/X)^o}
\operatorname{Hom}_{K(\mathcal A)}(\mathord\cdot,Y).
\]

\VerdierSemanticAnchor{ch3d-3-1-3-2}
\[
\textit{(2)}\qquad
\operatorname{Ext}_{\mathcal A}^{0}(X,Y)\simeq
\varinjlim_{Y\backslash\operatorname{Qis}}
\operatorname{Hom}_{K(\mathcal A)}(X,\mathord\cdot).
\]

\VerdierSemanticAnchor{ch3d-3-1-3-3}
\[
\textit{(3)}\qquad
\operatorname{Ext}_{\mathcal A}^{0}(X,Y)\simeq
\varinjlim_{(\operatorname{Qis}/X)^o\times
  (Y\backslash\operatorname{Qis})}
\operatorname{Hom}_{K(\mathcal A)}(\mathord\cdot,\mathord\cdot).
\]

\VerdierSemanticAnchor{ch3d-3-1-3-4}
\[
\textit{(4)}\qquad
\operatorname{Ext}_{\mathcal A}^{n}(X,Y)\simeq
\operatorname{Ext}_{\mathcal A}^{0}(X[p],Y[n+p]),
\qquad\textit{for every integer $p$.}
\]

\VerdierSemanticAnchor{ch3d-3-1-3-5}
\[
\textit{(5) For every $X$ in $K^{-}(\mathcal A)$,}\qquad
\operatorname{Ext}_{\mathcal A}^{0}(X,Y)\simeq
\varinjlim_{(\operatorname{Qis}^{-}/X)^o}
\operatorname{Hom}_{K(\mathcal A)}(\mathord\cdot,Y).
\]

\VerdierSemanticAnchor{ch3d-3-1-3-6}
\[
\textit{(6) For every $Y$ in $K^{+}(\mathcal A)$,}\qquad
\operatorname{Ext}_{\mathcal A}^{0}(X,Y)\simeq
\varinjlim_{Y\backslash\operatorname{Qis}^{+}}
\operatorname{Hom}_{K(\mathcal A)}(X,\mathord\cdot).
\]
\endgroup
